\subsubsection{Rotating calipers (polygon diameter)}

\paragraph{Idea intuitiva.}
Encuentra el \textbf{diametro} de un poligono convexo (par de vertices mas distantes) en $O(n)$. Asume poligono en orden counter-clockwise. Algoritmo: dos indices $i$ (sobre las aristas) y $k$ (sobre los vertices opuestos); en cada paso, rotar $k$ mientras el area del paralelogramo aumenta. La distancia $i$-$k$ puede ser maxima en cualquier momento del barrido.

\paragraph{Propiedades clave (invariantes).}
\begin{itemize}\setlength\itemsep{0pt}
	\item Cada indice avanza como maximo $n$ veces total.
	\item El ``diametro'' es la maxima distancia entre dos vertices.
	\item Retorna \textbf{distancia al cuadrado}.
\end{itemize}

\paragraph{Complejidad.}
\begin{itemize}\setlength\itemsep{0pt}
	\item $O(n)$.
\end{itemize}

\paragraph{Donde tocar para modificar.}
\begin{itemize}\setlength\itemsep{0pt}
	\item \textbf{Reconstruir el par de vertices}: aniadir indices.
	\item \textbf{Anchura minima / maxima}: variantes de calipers.
	\item \textbf{Distancia minima entre poligonos convexos}: otra variante.
\end{itemize}

\paragraph{Trampas y errores frecuentes.}
\begin{itemize}\setlength\itemsep{0pt}
	\item Asume poligono \textbf{estrictamente} convexo (sin vertices colineales). Si hay colineales, aniadir perturbacion o filtrar.
\end{itemize}

\paragraph{Codigo.}
Ver \texttt{cpp.tex}, seccion \textit{Rotating calipers (polygon diameter)}.
